\documentclass{article}
\usepackage[margin=25mm]{geometry}
\usepackage{amsmath,amssymb}
\usepackage{xurl}
\usepackage[hidelinks]{hyperref}
\setlength{\emergencystretch}{2em}
% Shared commands for notes in docs/paper-gaps/.

\DeclareUnicodeCharacter{2080}{\ifmmode{}_{0}\else\textsubscript{0}\fi}
\DeclareUnicodeCharacter{2081}{\ifmmode{}_{1}\else\textsubscript{1}\fi}
\DeclareUnicodeCharacter{2082}{\ifmmode{}_{2}\else\textsubscript{2}\fi}
\DeclareUnicodeCharacter{2099}{\ifmmode{}_{n}\else\textsubscript{n}\fi}

\newcommand{\Lean}{\textsc{Lean}}
\ExplSyntaxOn
\NewDocumentCommand{\leanid}{m}{
  \ifmmode
    \textsf{\detokenize{#1}}
  \else
    \group_begin:
    \urlstyle{sf}
    \tl_set:Nn \l_tmpa_tl {#1}
    \tl_replace_all:Nnn \l_tmpa_tl {\_} {_}
    \exp_args:NV \path \l_tmpa_tl
    \group_end:
  \fi
}
\ExplSyntaxOff
\providecommand{\tr}{\operatorname{tr}}
\newcommand{\tnleanIssueBase}{https://github.com/LionSR/TNLean/issues/}
\newcommand{\tnleanPrBase}{https://github.com/LionSR/TNLean/pull/}
\newcommand{\tnleanDocBase}{https://sirui-lu.com/TNLean/docs/}
\newcommand{\ghissue}[1]{\href{\tnleanIssueBase#1}{GitHub issue \##1}}
\newcommand{\ghpr}[1]{\href{\tnleanPrBase#1}{PR \##1}}
\newcommand{\doclink}[1]{\href{\tnleanDocBase#1.html}{\path{#1}}}

% Dirac notation and the identity operator, as in blueprint/src/macros/common.tex.
% \providecommand keeps note-local definitions authoritative.
\usepackage{dsfont}
\providecommand{\ket}[1]{|#1\rangle}
\providecommand{\bra}[1]{\langle #1|}
\providecommand{\braket}[2]{\langle #1 | #2 \rangle}
\providecommand{\ketbra}[2]{|#1\rangle\!\langle #2|}
\providecommand{\Id}{\mathds{1}}
\providecommand{\one}{\mathds{1}}
\providecommand{\C}{\mathbb{C}}
\providecommand{\MN}[1]{M_{#1}(\C)}
\providecommand{\MD}{M_D(\C)}

% Verdict marker: \gapnote{<kind>}{<status>}. Kind is the policy
% classification (clarification, local-correction, scope-restriction,
% unfaithful, false-source, open-gap); status is open, wip (elimination
% underway), resolved, or historical. Machine-read by the site index;
% prints a verdict line.
\newcommand{\gapnote}[2]{\par\noindent\textbf{Verdict:} #1 (#2).\par}

\title{Corrected regional soft-truncation residuals and parameter bounds}
\author{}
\date{2026-11-02}
\begin{document}
\maketitle
\gapnote{scope-restriction}{open}

\section{The source permits an estimated regional energy}

Definition \texttt{def:soft-trunc} of
\href{https://arxiv.org/abs/1602.08828v2}{ALVV, arXiv:1602.08828v2},
source lines 1178--1185, defines the regional replacement as
\begin{align}
 \widetilde A=e'I+h_{m,t}(A-e'I),\qquad e-10\leq e'\leq e,
\end{align}
where $e$ is the actual minimum energy of the regional Hamiltonian $A$.
Lemma \texttt{lem:exp-approx}, beginning at line 1195, allows $t\geq m\geq2$.
The ten-unit interval is part of the source statement; replacing $e'$ by
$e$ proves a specialization, not uniformity over all allowed estimates.

\section{The corrected bounds presently proved}

For positive-contraction nearest-neighbour interactions, the formalization
constructs the genuine minimum energies $e_0,e_1,e_2$ of three consecutive
internal regional sums $A_0,A_1,A_2$, together with normalized attaining
vectors. It centers the full original Hamiltonian by their sum:
\begin{align}
 H=\sum_i h_i-\Bigl(\sum_j e_j\Bigr)I,
 \qquad X=A_1-e_1I,
 \qquad G=H-X+h_{m,t}(X).
\end{align}
The attained original and soft ground energies $E$ and $\mu$ satisfy
$0\leq\mu\leq E\leq2$. Positivity of the actual regional sums and the
actual crossing bonds gives $G\geq0$; functional calculus gives $G\leq H$.
No gap or simplicity assumption is imposed.

Write $V_M(c)$ for the closed spectral space of $M$ below $c$, and put
$\varepsilon_{m,t}(T)=(t/m)(T/t)^m$. For $m\geq2$ and $t>0$, the same
constructed $G$ satisfies the following two statements:
\begin{align}
 \|(H-G)|_{V_H(E+\eta)}\|
 &\leq32e^{-(T-\eta)/8}+\varepsilon_{m,t}(T),
 &&\eta>0,\quad T\geq\eta+7,\\
 \|(H-G)|_{V_G(\mu+\eta)}\|
 &\leq16e^{-(T-U)/8}+\varepsilon_{m,t}(T),
 &&\eta\geq0,\quad U=2(\eta+8),\quad t,T\geq U.
\end{align}
The regional minima are selected before $t,m$; the actual soft operator and
its ground energy are selected before $\eta,T$.

The reverse proof obtains its coercive floor from
$h_{m,t}(x)\geq U-\varepsilon_{m,t}(U)$ for $x\geq U$.
Since $t\geq U$ and $m\geq2$,
$\varepsilon_{m,t}(U)\leq U/m\leq U/2$.
The floor is therefore at least $\eta+8$. The actual boundary contains at
most two bonds, so the Sylvester denominator is at least
$(\eta+8)-(\mu+\eta)-4\geq2$.
The boundary count, energy bound, commutation and projection inequality
are derived from the original interactions. They are not supplied as
hypotheses of the physical theorem.

The corresponding formal statements are
\leanid{TNLean.MPS.AreaLaw.ALVV.exists_single_region_original_soft_truncation_residual}
and
\leanid{TNLean.MPS.AreaLaw.ALVV.exists_single_region_soft_truncation_residuals}.

These are separate corrected auxiliary statements. They do not assert the
printed uniform energy estimate, its printed spectral enlargement, or the
complete efficient-filter construction. The separate counterexamples to the printed cluster estimate are recorded in
\href{https://sirui-lu.com/TNLean/paper-gaps/alvv16_cluster_approximation.pdf}{the cluster-approximation note}.

\section{Estimated centers and simultaneous replacements}

The energy-center restriction is removed in a further corrected auxiliary
statement. Four ordered cuts determine five consecutive regions with
internal sums $A_j$ and an actual crossing set $J$. Put $C=|J|\leq4$.
The original interactions determine attained regional minima $e_j$ and
an attained original ground energy $E\in[0,C]$ before every selected set
$S$, every family of estimates $p_j\in[e_j-10,e_j]$ on $S$, and every
scale $t>0$ and order $m\geq2$. Define
\begin{align}
 B&=\sum_j\left(
  \begin{cases}
   p_jI+h_{m,t}(A_j-p_jI),&j\in S,\\
   A_j,&j\notin S
  \end{cases}-e_jI\right),\\
 D&=10|S|,\qquad G=B+\sum_{i\in J}h_i.
\end{align}
The actual common interior satisfies $B+DI\geq0$ and commutes with
every $A_j$. The attained soft ground energy obeys
$-D\leq\mu\leq E\leq C$. For $\eta\geq0$,
\begin{align}
 \|(H-G)|_{V_H(E+\eta)}\|
 &\leq |S|\left(8Ce^{-(T-\eta-3C-12)/8}+\varepsilon_{m,t}(T)\right),
 &&T\geq\eta+3C+12,\\
 \|(H-G)|_{V_G(\mu+\eta)}\|
 &\leq |S|\left(8Ce^{-(T-U)/8}+\varepsilon_{m,t}(T)\right),
 &&U=2(\eta+D+3C+2),\quad t,T\geq U.
\end{align}
The same $G$, ground data and actual orthonormal eigenbases precede
$\eta,T$. The boundary count, positivity correction, commutation and
projection floors are derived from the original local terms.

In the reverse argument the shifted soft energy is at most $C+D+\eta$,
while the high-energy floor is at least $U/2=\eta+D+3C+2$.
Consequently the actual denominator is at least two after subtracting
the boundary contribution $2C$. No positivity or spectral gap of $G$
is assumed. If the source's three bulk regions are selected, then
$|S|\leq3$, $D\leq30$ and $U\leq2(\eta+44)$.

The public physical statement is
\leanid{TNLean.MPS.AreaLaw.ALVV.exists_five_region_soft_residuals}.
The accompanying regional construction is
\leanid{TNLean.MPS.AreaLaw.ALVV.exists_five_region_soft_structure}.
The two restricted-norm estimates for a general common regional
interior are
\leanid{TNLean.MPS.AreaLaw.ALVV.norm_simultaneous_soft_original_residual}
and
\leanid{TNLean.MPS.AreaLaw.ALVV.norm_simultaneous_soft_reverse_residual}.
Empty intervals, repeated cuts and an empty selected set are included.

\section{The remaining source comparison}

The exact-center specialization no longer prevents a simultaneous
estimated-center comparison. The open qualification concerns the
explicit scale and cutoff conditions and the subsequent conclusions:
the corrected residuals must be composed with the actual spectral
comparison, the same represented soft Hamiltonian and the source's
filter parameter choices. They do not by themselves assert the printed
uniform energy estimate, its spectral enlargement, a polynomial
construction with the source's dependence on varying parameters, or a
running-time theorem. In particular, a bound polynomial in the chain
length for fixed order cannot substitute for a bound uniform in an
order chosen as a function of that length.

The source's GEN input at lines 1360--1362 permits a symmetric estimate
$|\varepsilon'_M-\varepsilon_M|\leq10$, whereas Definition
\texttt{def:soft-trunc} uses a lower estimate in a ten-unit interval.
The present statement has exactly the latter one-sided range. Passing
from the GEN input to suitable actual regional lower centers remains
part of the algorithmic argument; it is not inferred from the existence
of the true minimum energies. The separate false cluster estimate and
its counterexamples are unchanged.
\end{document}
